% Options for packages loaded elsewhere
\PassOptionsToPackage{unicode}{hyperref}
\PassOptionsToPackage{hyphens}{url}
\documentclass[
  11pt,
]{article}
\usepackage{xcolor}
\usepackage[margin=25mm]{geometry}
\usepackage{amsmath,amssymb}
\setcounter{secnumdepth}{-\maxdimen} % remove section numbering
\usepackage{iftex}
\ifPDFTeX
  \usepackage[T1]{fontenc}
  \usepackage[utf8]{inputenc}
  \usepackage{textcomp} % provide euro and other symbols
\else % if luatex or xetex
  \usepackage{unicode-math} % this also loads fontspec
  \defaultfontfeatures{Scale=MatchLowercase}
  \defaultfontfeatures[\rmfamily]{Ligatures=TeX,Scale=1}
\fi
\usepackage{lmodern}
\ifPDFTeX\else
  % xetex/luatex font selection
\fi
% Use upquote if available, for straight quotes in verbatim environments
\IfFileExists{upquote.sty}{\usepackage{upquote}}{}
\IfFileExists{microtype.sty}{% use microtype if available
  \usepackage[]{microtype}
  \UseMicrotypeSet[protrusion]{basicmath} % disable protrusion for tt fonts
}{}
\makeatletter
\@ifundefined{KOMAClassName}{% if non-KOMA class
  \IfFileExists{parskip.sty}{%
    \usepackage{parskip}
  }{% else
    \setlength{\parindent}{0pt}
    \setlength{\parskip}{6pt plus 2pt minus 1pt}}
}{% if KOMA class
  \KOMAoptions{parskip=half}}
\makeatother
\setlength{\emergencystretch}{3em} % prevent overfull lines
\providecommand{\tightlist}{%
  \setlength{\itemsep}{0pt}\setlength{\parskip}{0pt}}
\usepackage{bookmark}
\IfFileExists{xurl.sty}{\usepackage{xurl}}{} % add URL line breaks if available
\urlstyle{same}
\hypersetup{
  hidelinks,
  pdfcreator={LaTeX via pandoc}}

\author{}
\date{}

\begin{document}

\section{The index pairing between K-theory and
K-homology}\label{the-index-pairing-between-k-theory-and-k-homology}

\emph{Self-checked by the writing AI, GPT-6.1 Sol (OpenAI). Public
domain (CC0).}

A Fredholm module supplies an operator that almost intertwines the
algebra action. A projection selects two subspaces on which that
operator is Fredholm; a unitary supplies a Fredholm compression on the
positive spectral subspace. Their finite-dimensional defects are the
even and odd index pairings.

Our convention is \[
\operatorname{index}T=\dim\ker T-\dim\ker T^*.
\tag{0.1}
\] For an odd exact involution \(F\), the projection is \(P=(1+F)/2\).
With these two choices the Hardy module pairs with the circle coordinate
\(z\) to \(-1\). Section 7 proves its trace-character expression with
the same positive-compression convention.

We use the preceding lesson's normalization, group laws and bounded
transform, the Hilbert-space Fredholm facts specified there, and the
projection and unitary definitions of K-theory from the operator
K-theory course. For the boundary map we use its published lesson
\emph{The index map and the exact sequence at \(K_0\)}, Definition 2.1
and Theorem 2.2. We compute its sign here from a doubled lift.

\subsection{1. The even compression}\label{the-even-compression}

Let \(A\) first be unital and let \(x=(H,\pi,F)\) be an even cycle.
Matrix amplification means \(H^n=H\otimes\mathbb C^n\), with grading
\(\Gamma\otimes1\), operator \(F^{(n)}=F\otimes1\), and matrix
representation \(\pi^{(n)}([a_{ij}])=[\pi(a_{ij})]\). For a projection
\(p\in M_n(A)\), put \[
\begin{aligned}
P_\pm&=\pi_\pm^{(n)}(p),\\
T_p&=P_-F_+^{(n)}|_{P_+H_+^n}.
\end{aligned}
\tag{1.1}
\] The domain and codomain in this formula are \(P_+H_+^n\) and
\(P_-H_-^n\); they can have different dimensions.

\subsubsection{Lemma 1.1. The even compression is
Fredholm}\label{lemma-1.1.-the-even-compression-is-fredholm}

\(T_p\) is Fredholm, with inverse modulo compacts
\(S_p=P_+F_-^{(n)}|_{P_-H_-^n}\). Its index is unchanged by the
preceding lesson's normalization.

\textbf{Proof.} Suppress the amplification notation. The off-diagonal
blocks of \([F,\pi(p)]\) give \[
\begin{aligned}
F_+P_+-P_-F_+&\in\mathcal K,\\
F_-P_--P_+F_-&\in\mathcal K.
\end{aligned}
\] The two inverse errors have the exact typed expansions \[
\begin{aligned}
S_pT_p-P_+
&=P_+F_-(P_-F_+-F_+P_+)P_+
  +P_+(F_-F_+-1)P_+,\\
T_pS_p-P_-
&=P_-F_+(P_+F_--F_-P_-)P_-
  +P_-(F_+F_--1)P_- .
\end{aligned}
\] The first terms are compact by the preceding commutators; the second
terms are the represented corners of the square defect in (1.1). The
adjoint defect also gives \(T_p^*-S_p=P_+(F_+^*-F_-)P_-\in\mathcal K\).
All these compact operators have the indicated represented ranges as
their domain and target. The Fredholm criterion in the first extension
lesson therefore applies even when the two ranges have different
dimensions or the representation is degenerate.

Replacing \(F\) by its self-adjoint clipped contraction changes its
product with \(\pi(p)\) by a compact. Thus the compressed block changes
compactly. In the doubled exact normalization the representation is
\(\pi\oplus0\), so the projection \(p\) selects only the first
represented copy. Its compressed operator is exactly that clipped block;
the off-diagonal square-root entries are killed by the zero
representation. Index invariance under compact perturbation proves the
last assertion. \(\square\)

Define \(\langle[p],x\rangle=\operatorname{index}T_p\), and extend to
differences of projections by subtraction.

\subsubsection{Theorem 1.2. The even pairing for unital
algebras}\label{theorem-1.2.-the-even-pairing-for-unital-algebras}

This defines a bilinear pairing \[
K_0(A)\times K^0(A)\longrightarrow\mathbb Z.
\tag{1.2}
\] It is invariant under projection equivalence and homotopy, operator
homotopy, compact perturbations, and addition of degenerate cycles.

\textbf{Proof.} A block sum \(p\oplus q\) gives \(T_p\oplus T_q\), so
the index is additive in projections and insensitive to zero padding. If
\(v\) is a partial isometry implementing \(v^*v=p\), \(vv^*=q\), its two
represented blocks give unitaries
\(W_\pm:P_\pm H_\pm^n\to Q_\pm H_\pm^m\). Rectangular matrix entries of
the cycle commutator give \[
W_-T_p-T_qW_+\in\mathcal K.
\] In detail that difference is
\(Q_-(\pi_-^{(m,n)}(v)F_+^{(n)}-F_+^{(m)}\pi_+^{(m,n)}(v))P_+\), whose
finitely many matrix entries are the cycle commutators. The two
compressions therefore agree up to compact error after these unitary
identifications, and have equal index. This proves independence under
stable projection equivalence and hence extension to the Grothendieck
group.

For a norm-continuous path of projections, their represented ranges are
locally identified by the continuous close-projection unitaries proved
in the preceding lesson, Lemma 7.1. Use those unitaries separately in
the two graded spaces. The transported compressions are a
norm-continuous Fredholm path, so the index is constant. An operator
homotopy of \(F\) gives such a path on the fixed represented projection
ranges directly. A locally compact change of \(F\) is compact on those
ranges, and a compact perturbation is a special case.

Unitary equivalence of cycles conjugates the typed operator. A direct
sum of cycles gives a direct sum of typed operators, so the index is
additive in cycles. For a degenerate cycle, the two compressed blocks
are actual inverse adjoints: the commutator, adjoint defect and square
defect are all zero on the represented ranges. Thus \(T_p\) is unitary
from one range to the other and its index is zero. Every generating
K-homology relation preserves the index. Additivity in each variable
proves bilinearity. \(\square\)

For a unital representation, testing \(p=1\) gives
\(\operatorname{index}F_+\). A finite-dimensional scalar cycle with
graded dimensions \(r,s\) therefore gives \(r-s\), as announced in the
preceding lesson.

\subsection{2. The odd compression}\label{the-odd-compression}

First take an odd cycle with \(F=F^*\), \(F^2=1\), and unital
representation. Put \(P=(1+F)/2\). For \(u\in M_n(A)\) unitary define \[
T_u=P^{(n)}\pi^{(n)}(u)|_{P^{(n)}H^n},
\qquad
\langle[u],x\rangle=\operatorname{index}T_u.
\tag{2.1}
\]

\subsubsection{Lemma 2.1. Fredholmness and the product
rule}\label{lemma-2.1.-fredholmness-and-the-product-rule}

The operator \(T_u\) is essentially unitary, and \[
T_{uv}-T_uT_v\in\mathcal K,\qquad
\operatorname{index}T_{uv}
=\operatorname{index}T_u+\operatorname{index}T_v.
\tag{2.2}
\]

\textbf{Proof.} Write \(U=\pi^{(n)}(u)\). Since \([P^{(n)},U]\) is
compact, the differences \[
T_u^*T_u-1_{PH^n}
=-P^{(n)}U^*(1-P^{(n)})U|_{PH^n}
\] and \(T_uT_u^*-1_{PH^n}\) are compact. Thus \(T_{u^*}=T_u^*\) is an
inverse modulo compacts. For another represented unitary \(V\), \[
T_{uv}-T_uT_v
=P^{(n)}U(1-P^{(n)})V|_{PH^n}
\] is compact. Compact invariance and Fredholm composition additivity
give (2.2). \(\square\)

\subsubsection{Theorem 2.2. The odd pairing for unital
algebras}\label{theorem-2.2.-the-odd-pairing-for-unital-algebras}

Formula (2.1) extends to a bilinear pairing \[
K_1(A)\times K^1(A)\longrightarrow\mathbb Z,
\tag{2.3}
\] using the continuous exact normalization for a general cycle. It is
invariant under all K-theory and K-homology relations.

\textbf{Proof.} For a possibly degenerate representation, replace the
represented unitary by \(U=\pi^{(n)}(u)+1-\pi^{(n)}(1)\). This is an
actual unitary and acts as the identity on the representation's zero
part. For a general \(F\), use the exact doubled normalization of the
preceding lesson, Proposition 2.3, and this same unitary completion. Its
Lemma 7.1, applied at each matrix size, proves invariance under operator
homotopy, unitary equivalence, degeneracy and direct sum of cycles, and
proves that removing the zero-representation part contributes zero
index.

A unitary path \(u_t\) gives a norm-continuous Fredholm compression path
on the fixed \(PH^n\). Its index is constant. Matrix padding by an
identity adds an identity compression of index zero; block sums add the
indices. These are the stable-unitary relations defining \(K_1(A)\). The
product rule is consistent with that group law: the path \[
\operatorname{diag}(u,1)R_t\operatorname{diag}(1,v)R_t^*,
\quad
R_t=\begin{pmatrix}\cos t&-\sin t\\
                         \sin t&\cos t\end{pmatrix},
\] for \(0\leq t\leq\pi/2\), joins \(\operatorname{diag}(u,v)\) to
\(\operatorname{diag}(uv,1)\). Thus addition may be computed by either
block sum or product. Additivity in the cycle and unitary variables
proves the theorem. Compact perturbations of \(F\) give operator
homotopies by the preceding lesson, Proposition 2.2, so they preserve
the pairing. \(\square\)

Replacing \(u\) by \(u^*\) negates the index. Replacing \(F\) by \(-F\)
negates the cycle and hence the pairing. These are two different
operations with the same sign effect.

\subsection{3. Nonunital algebras and the scalar
part}\label{nonunital-algebras-and-the-scalar-part}

Write \(A^+=A\oplus\mathbb C\) for the forced unitization, adjoining a
unit even when \(A\) already has one. Let \(\epsilon:A^+\to\mathbb C\)
be its scalar quotient. The existing
\href{https://kokunoyumeto.github.io/open-math-courses/courses/KT-OPK/nonunital-algebras-unitization-relative-classes-and-half-exactness.html}{Nonunital
algebras: unitization, relative classes and half-exactness}, Section 1
and Theorem 1.1, gives \[
K_0(A)=\ker\epsilon_*,
\tag{3.1}
\] and represents a class there as \([p]-[q]\) for projections over
\(A^+\) whose scalar ranks agree. The existing
\href{https://kokunoyumeto.github.io/open-math-courses/courses/KT-OPK/invertibles-unitaries-and-k1.html}{Invertibles,
unitaries and \(K_1\)}, Definition 1.1 and Proposition 1.2, gives stable
unitary representatives over \(A^+\); their scalar part may be
normalized to the identity.

A raw nonunital cycle need not extend to a cycle over \(A^+\), because
its global square defect may fail to be compact. This is the reason for
normalizing before extending.

\subsubsection{Theorem 3.1. Relative
pairings}\label{theorem-3.1.-relative-pairings}

For every separable \(A\), the preceding pairings have natural
extensions \[
K_i(A)\times K^i(A)\longrightarrow\mathbb Z,\qquad i=0,1.
\tag{3.2}
\] For an exact normalized cycle they are the unitization compression
formulas, with subtraction for a relative projection class. Scalar
projection and unitary classes contribute zero.

\textbf{Proof.} Given any cycle, perform the particular exact
normalization \(N(F)\) of the preceding lesson on
\(\widetilde H=H\oplus H\), using the opposite second grading in the
even case and representation \(\rho=\pi\oplus0\). Define the unital
representation \[
\rho^+(a+\lambda1)=\rho(a)+\lambda1_{\widetilde H}.
\tag{3.3}
\] The normalized operator is an exact self-adjoint involution and its
commutators with \(\rho(A)\) are compact; scalar multiplication commutes
with it. Thus it is a cycle over \(A^+\).

In the even case let \(E_x(p)\) be the even compression index for
\(\rho^+(p)\), and define \[
\langle[p]-[q],x\rangle=E_x(p)-E_x(q).
\tag{3.4}
\] Theorem 1.2 shows that \(E_x\) is a homomorphism on \(K_0(A^+)\), so
its restriction to (3.1) is independent of the particular equal-rank
representatives. For a scalar projection \(p_0\in M_n(\mathbb C)\), its
representation commutes exactly with \(N(F)\). The positive-to-negative
block of this exact odd involution is a unitary, and restricts to a
unitary on the two scalar-selected graded ranges. Hence \(E_x(p_0)=0\).

In the odd case define the index using \(\rho^{+,(n)}(u)\), for
\(u\in M_n(A^+)\), and the positive projection of \(N(F)^{(n)}\).
Theorem 2.2 proves independence of stable-unitary representatives. A
scalar unitary acts only on the matrix factor, commutes exactly with
that projection and has a unitary compression, so its index is zero.
Multiplying \(u\) by \(\epsilon(u)^*\) therefore changes neither its
K-theory class nor its index. The scalar matrix has a path to the
identity, by finite-dimensional spectral diagonalization and
exponentiation, justifying the normalized representative convention.

It remains to check dependence only on the original \(A\)-cycle class,
since unitization of an arbitrary homotopy was not available before
normalization. The construction \(F\mapsto N(F)\) is continuous. Thus an
operator homotopy over \(A\) becomes one over \(A^+\): its square and
adjoint defects are exactly zero, and its compact commutators are those
over \(A\). Unitary equivalence and direct sum commute with the
construction up to permutations. For a degenerate original cycle, the
proof in the preceding lesson, Lemma 7.1, shows that its normalized
operator commutes exactly with \(\rho(A)\), hence also with (3.3). The
normalized cycle over \(A^+\) is then degenerate, so both its pairings
vanish. This verifies every generating relation.

If an exact representative was already used, the further normalization
adds a zero \(A\)-representation on which the operator is \(-F\). On
extending to \(A^+\), this second summand has only scalar action and an
exact involution; all its compression indices are zero by the preceding
scalar checks. Thus (3.4) and the odd formula can be evaluated directly
on any exact representative, and agree with the canonical construction.
For projections already in \(A\), Lemma 1.1 recovers the original even
compression. For a unital algebra and \(u\in M_n(A)\), its normalized
forced-unitization representative is \(u+1-1_A\), whose representation
is \(\rho(u)+1-\rho(1_A)\). This recovers precisely the earlier unital
formulas. \(\square\)

For a *-homomorphism \(\alpha:A\to B\), these pairings satisfy \[
\langle\alpha_*k,x\rangle=\langle k,\alpha^*x\rangle.
\tag{3.5}
\] Indeed \(\alpha\) extends unitally to the forced unitizations, and
the amplified represented projections or unitaries on both sides are
identical. The normalized operator depends on \(F\), which is also
identical. This proves naturality by the compression formulas.

\subsection{4. The extension boundary and its
sign}\label{the-extension-boundary-and-its-sign}

Let an exact odd cycle give the compression Busby map
\(\tau:A\to\mathcal Q(PH)\). Its pullback extension is semisplit. After
the infinite zero-representation padding described in the preceding
lesson, we may assume that \(PH\) is infinite dimensional.

\subsubsection{Lemma 4.0. The rank identification for the compact
ideal}\label{lemma-4.0.-the-rank-identification-for-the-compact-ideal}

For an infinite-dimensional Hilbert space \(L\), finite rank gives an
isomorphism \(K_0(\mathcal K(L))\cong\mathbb Z\), taking a rank-one
projection to \(1\). Matrix amplification preserves this identification.
In particular the class \([p_k]-[p_c]\) of finite-rank projections has
value \(\operatorname{rank}p_k-\operatorname{rank}p_c\).

\textbf{Proof.} A compact projection has finite-dimensional range: an
infinite orthonormal family in its range would contradict compactness of
its image of the unit ball. Two finite-rank projections of equal rank
are Murray--von Neumann equivalent by any isometry between their ranges,
extended by zero. The implementing operator is finite rank. Thus their
classes equal the appropriate multiple of a fixed rank-one class.

We must also account for all relative projections and show that no extra
relation kills that class. Represent \(M_n(\mathcal K(L)^+)\) on
\(L^n\), and for a projection \(p\) denote its scalar projection by
\(p_0\). The difference \(p-p_0\) is compact. The map \[
A_p=pp_0:p_0L^n\longrightarrow pL^n
\] is Fredholm: \(p_0p\) is an inverse modulo compact operators, since
the errors are products with \(p-p_0\). Define
\(r(p)=-\operatorname{index}A_p\). This integer is additive under block
sums. It is unchanged under stable projection equivalence: if
\(v^*v=p,\ vv^*=q\), its scalar part \(v_0\) satisfies
\(v_0^*v_0=p_0,\ v_0v_0^*=q_0\). Hence \(v\) and \(v_0\) are unitaries
between the corresponding represented ranges, and \[
vA_p-A_qv_0=q(v-v_0)p_0
\] is compact. Unitary identifications and compact invariance of the
index give \(r(p)=r(q)\). This also treats rectangular matrices after
padding. Along a projection path, both \(p\) and \(p_0\) vary in norm;
the close-projection unitaries from the preceding lesson transport both
ranges locally to fixed spaces, so index invariance makes \(r\)
constant. Consequently \(r\) descends to the projection Grothendieck
group. It is zero on scalar projections and equals the rank on
finite-rank projections, whose scalar part is zero.

Direct all finite-dimensional subspaces of \(L\) by inclusion and denote
their projections by \(e_\alpha\). This net converges strongly to
\(1_L\), since each vector belongs to one of those subspaces. Put
\(E_\alpha=1_n\otimes e_\alpha\); these commute with \(p_0\). A
finite-rank approximation followed by convergence on its
finite-dimensional domain and range proves \[
\|E_\alpha(p-p_0)E_\alpha-(p-p_0)\|\longrightarrow0.
\] Thus the self-adjoint \(b_\alpha=p_0+E_\alpha(p-p_0)E_\alpha\)
converges to \(p\). Eventually its spectrum is confined to disjoint
neighborhoods of \(0,1\); this follows from the Neumann series for
\(b_\alpha-\lambda\) whenever the distance from \(\lambda\) to
\(\{0,1\}\) exceeds \(\|b_\alpha-p\|\). A continuous spectral cutoff
gives a projection \(q_\alpha\) with \(\|q_\alpha-p\|\to0\). Continuity
here follows by uniformly approximating that fixed cutoff by
polynomials. The construction respects the blocks \(E_\alpha L^n\) and
its complement; on the complement \(q_\alpha=p_0\).

Choose such an \(\alpha\) with \(\|q_\alpha-p\|<1\). These close
projections are conjugate by the explicit unitary \(X(X^*X)^{-1/2}\),
where \(X=q_\alpha p+(1-q_\alpha)(1-p)\). It belongs to
\(M_n(\mathcal K(L)^+)\): \(q_\alpha-p\) is compact, so \(X-1\) and, by
quotient functional calculus, \((X^*X)^{-1/2}-1\) are compact. Write the
finite-dimensional restriction of \(q_\alpha\) as \(q_f\). For
orthogonal projections \(a,b\), the row \((a,b)\) has products \(a+b\)
and \(\operatorname{diag}(a,b)\), so implements their stable equivalence
and proves direct-sum addition. Applying this to the two blocks yields
\[
[p]-[p_0]=[q_\alpha]-[p_0]=[q_f]-[p_0E_\alpha].
\] Both projections on the right are finite rank.

Every relative class in \(K_0(\mathcal K(L))\), defined as the kernel of
the scalar quotient on the forced unitization, is \([p]-[q]\) with equal
scalar ranks. The finite-dimensional scalar projections \(p_0,q_0\) are
stably equivalent, so the preceding reduction writes this class as a
difference of finite-rank projection classes. All classes are therefore
integer multiples of the rank-one class. The integer \(r\) constructed
above sends that class to \(1\), so there is no additional relation.
This proves the isomorphism. All arguments on \(L^n\) use ordinary
finite rank and give exactly the same integer after amplification.
\(\square\)

\subsubsection{Theorem 4.1. The odd pairing is the index
boundary}\label{theorem-4.1.-the-odd-pairing-is-the-index-boundary}

For this extension, \[
\partial:K_1(A)\longrightarrow K_0(\mathcal K(PH))\cong\mathbb Z
\] is exactly \(k\mapsto\langle k,x\rangle\), when the rank-one compact
projection is the positive generator of \(K_0(\mathcal K)\).

\textbf{Proof.} Use a normalized unitary \(u\) over \(A^+\). Its
compression \(T\) is a lift of the unitary \(\tau^{+,(n)}(u)\) in the
Calkin algebra. By Lemma 2.1 it is essentially unitary and Fredholm. Its
polar partial isometry \(V\) has \[
V^*V=1-p_k,\qquad VV^*=1-p_c,
\tag{4.1}
\] where \(p_k,p_c\) are the finite-rank projections onto its kernel and
cokernel. For an explicit construction, Fredholm closed range gives a
gap in the spectrum of \(|T|\) off zero. The continuous function \(h\)
on that spectrum, with \(h(0)=0\) and \(h(t)=1/t\) for \(t>0\), gives
\(V=Th(|T|)\) and the identities (4.1). Also \(T-V\) is compact. To
check this, \(|T|^2-1\) is compact, so \(|T|-1\) is compact by
continuous functional calculus in the Calkin quotient. The identity
\(T=V|T|\) gives \(T-V=V(|T|-1)\).

There is an explicit doubled unitary lift: \[
W=\begin{pmatrix}V&-p_c\\p_k&V^*\end{pmatrix}.
\tag{4.2}
\] The relations \(Vp_k=0\), \(p_cV=0\), and their adjoints show by
multiplication that \(WW^*=W^*W=1\). Its quotient has diagonal entries
\(\tau^{+,(n)}(u)\) and its adjoint, and zero off-diagonal entries. With
\(P_0=\operatorname{diag}(1,0)\), \[
WP_0W^*=\operatorname{diag}(1-p_c,p_k).
\tag{4.3}
\] The boundary definition in the operator K-theory prerequisite is the
difference of this projection and \(P_0\), in the compact ideal's
relative \(K_0\). It therefore gives \[
\partial[u]=[p_k]-[p_c]
\quad\longmapsto\quad
\dim\ker T-\dim\ker T^*.
\tag{4.4}
\] This is the pairing by definition.

The computation is compatible with the forced-unitization convention
even though the Calkin extension is already unital. In the forced
unitization of \(\mathcal B(PH^n)\), use the normalized doubled lift
\(w=W+(1_{\mathrm{ext}}-1_{\mathcal B})1_2\) and the reference
projection
\(P_0^{\mathrm{ext}}=\operatorname{diag}(1_{\mathrm{ext}},0)\). The two
orthogonal central summands \(1_{\mathcal B}\) and
\(1_{\mathrm{ext}}-1_{\mathcal B}\) make their products split. Thus \[
wP_0^{\mathrm{ext}}w^*-P_0^{\mathrm{ext}}
=WP_0W^*-P_0=\operatorname{diag}(-p_c,p_k).
\] In the compact ideal's unitization the resulting projection is
\(\operatorname{diag}(1-p_c,p_k)\), with reference
\(\operatorname{diag}(1,0)\). Orthogonal-sum addition, proved in Lemma
4.0, makes its relative class exactly \([p_k]-[p_c]\). This verifies the
scalar convention directly rather than identifying the two units
silently.

Formally, this Calkin boundary is the boundary of the pullback extension
by naturality: its projection to the multiplier algebra is the identity
on the compact ideal and induces the Busby map on the quotient.
\href{https://kokunoyumeto.github.io/open-math-courses/courses/KT-OPK/the-index-map-and-the-exact-sequence-at-k0.html}{The
index map and the exact sequence at \(K_0\)}, Definition 2.1 and Theorem
2.2, proves that naturality and the doubled-lift formula, including
choice of lift, normalized homotopy and block-sum addition. Lemma 4.0
proves the rank identification and its invariance under matrix
amplification here. Scalar unitary parts give zero, by Theorem 3.1, so
the argument proves the assertion for every class in \(K_1(A)\).
\(\square\)

In particular, for the Toeplitz extension the lift is \(S\), so
\(p_k=0\), \(p_c=p_0\), and \(\partial[z]=-[p_0]\). The boundary sign
agrees with \(-1\) in our Fredholm convention; no change of orientation
is hidden in the identification with \(\mathbb Z\).

\subsection{5. Why a bundle test is a twisted Dirac
index}\label{why-a-bundle-test-is-a-twisted-dirac-index}

An index pairing with a bundle is an analytic assertion before it is a
characteristic-class formula. We prove the analytic comparison here.

\subsubsection{Lemma 5.1. A bounded perturbation changes the transform
compactly}\label{lemma-5.1.-a-bounded-perturbation-changes-the-transform-compactly}

Suppose \(D_0,D_1\) are self-adjoint with the same domain and compact
resolvents, and \(D_1-D_0=B\) is bounded. Their bounded transforms
differ by a compact operator.

\textbf{Proof.} Use the resolvent integral (4.5) of the preceding
lesson, with \(\alpha(s)=\sqrt{1+s^2}\). The resolvent identity gives \[
(D_1\pm i\alpha)^{-1}-(D_0\pm i\alpha)^{-1}
=-(D_1\pm i\alpha)^{-1}B(D_0\pm i\alpha)^{-1}.
\] Each difference is compact and has norm at most \(\|B\|/(1+s^2)\).
The integral of the two differences therefore converges in norm to a
compact. The individual truncated transform integrals converge strongly;
their difference must have this compact norm limit. \(\square\)

\subsubsection{Lemma 5.1a. The domain under a bounded self-adjoint
perturbation}\label{lemma-5.1a.-the-domain-under-a-bounded-self-adjoint-perturbation}

If \(D\) is self-adjoint and \(B=B^*\) is bounded, then \(D+B\), with
domain \(\operatorname{Dom}D\), is self-adjoint. Its graph norm is
equivalent to that of \(D\). If \(D\) has compact resolvent, so does
\(D+B\). Oddness for a grading is preserved when both \(D\) and \(B\)
are odd.

\textbf{Proof.} A vector \(v\) is in the adjoint domain of \(D+B\)
exactly when \(u\mapsto\langle (D+B)u,v\rangle\) is bounded on
\(\operatorname{Dom}D\) in the ambient Hilbert norm. The \(B\)-term is
already bounded, so this is exactly the adjoint-domain condition for
\(D\). Self-adjointness of \(D,B\) proves the assertion and the equality
of actions. The two triangle inequalities
\(\|(D+B)u\|\leq\|Du\|+\|B\|\|u\|\) and
\(\|Du\|\leq\|(D+B)u\|+\|B\|\|u\|\) give graph-norm equivalence. For
\(t>\|B\|\), \[
D+B-it=(1+B(D-it)^{-1})(D-it).
\] The first factor is invertible by its norm-convergent Neumann series,
since \(\|(D-it)^{-1}\|\leq1/t\). Hence \((D+B-it)^{-1}\) is compact.
The resolvent identity expresses the resolvent at any other nonreal
point as this compact operator plus its product with bounded resolvents,
so all such resolvents are compact. Their product gives compact
\((1+(D+B)^2)^{-1}\). The grading assertions are identities on the
common domain. \(\square\)

\subsubsection{Theorem 5.2. The twisted Dirac
comparison}\label{theorem-5.2.-the-twisted-dirac-comparison}

Let \(D\) be the graded spin\(^{c}\) Dirac operator on a closed
even-dimensional manifold \(M\), and let \(E\) be a smooth Hermitian
vector bundle with a Hermitian connection. With \(x=[D]\) the
bounded-transform class, \[
\langle[E],x\rangle=\operatorname{index}D_E^+.
\tag{5.1}
\] The operator on the right has domain \(H^1(M,S^+\otimes E)\) and
target \(L^2(M,S^-\otimes E)\). No topological index formula is assumed.

\textbf{Proof.} The spinor Dirac operator has the local expression
\(D=\sum_j c(e_j)\nabla^S_{e_j}\), where \(c(e_j)^*=-c(e_j)\),
\(c(\xi)^2=-|\xi|^2\), and the spinor connection is Hermitian and
compatible with Clifford multiplication. Its principal symbol is
\(ic(\xi)\), whose square is \(|\xi|^2\), proving ellipticity.
Integration by parts on the closed manifold makes its formal adjoint
equal to itself: the connection differentiates \(c(e_j)\) as
\(c(\nabla_{e_j}e_j)\), and the zero-order correction is
\(c(\sum_j\nabla_{e_j}e_j+\sum_j\operatorname{div}(e_j)e_j)=0\), since
metric compatibility gives
\(\sum_j\operatorname{div}(e_j)e_j=-\sum_j\nabla_{e_j}e_j\). There is no
boundary term. Clifford multiplication is odd and the connection
preserves the grading. The preceding lesson, Corollary 4.2, with the
complete finite-chart and Friedrichs-mollifier proofs of Lemmas
4.1a--4.1c, therefore gives the self-adjoint \(H^1\) realization with
compact resolvent used below.

Embed \(E\) smoothly and isometrically in a trivial finite-rank bundle.
One direct construction chooses finitely many local orthonormal frames
and smooth real cutoffs \(\chi_j\) supported in their charts, with
\(\sum_j\chi_j^2=1\). Send a vector to the direct sum of its frame
coordinates multiplied by the \(\chi_j\), extending each product by
zero. This is a smooth isometry, so its image has smooth orthogonal
projection \(p\in M_N(C^\infty(M))\).

On \(L^2(M,S)\otimes\mathbb C^N\), with the amplified \(D\),
multiplication by \(p\) is even and preserves \(H^1\). Its commutator
with \(D\) is bounded by the first-order formula in the preceding
lesson. Put \[
D^\mathrm{diag}=pDp+(1-p)D(1-p).
\] This has the same domain as \(D\) and differs by the bounded
off-diagonal part \(pD(1-p)+(1-p)Dp\). The latter is bounded because
\(pD(1-p)=-p[D,p]\), and the other term is its adjoint. The off-diagonal
part is self-adjoint because its two blocks are adjoints, and it is odd
because \(p\) is even. Lemma 5.1a therefore gives self-adjointness of
\(D^{\mathrm{diag}}\) on the same \(H^1\) domain and its compact
resolvent, without requiring a second elliptic-domain assertion.

Lemma 5.1 shows that its bounded transform differs compactly from that
of \(D\). Since \(p\) preserves the domain and commutes with the
diagonal operator there, it commutes with each resolvent: apply
\((D^{\mathrm{diag}}\pm i)^{-1}\) to the domain identity. The
restrictions to \(pH\) and its complement therefore have surjective
nonreal resolvents and are self-adjoint, with compact resolvent. The
resolvent integral of the preceding lesson, Theorem 4.1, restricts on
\(pH\) to the same integral for \(D_p=pDp\), so the restricted transform
is precisely its bounded transform on \(S\otimes E\). Hence the even
compression defining \(\langle[p],[D]\rangle\) has the same index as
\(F_{D_p}^+\).

The operator \(D_p\) is the Dirac operator for the projected connection
on \(E\). Replacing that connection by the given Hermitian connection
changes the Dirac operator by a smooth bundle endomorphism, hence a
bounded self-adjoint odd operator. Explicitly the connection difference
is an endomorphism-valued one-form \(B_j\), with \(B_j^*=-B_j\) because
both connections preserve the Hermitian metric. The operator difference
is \(\sum_j c(e_j)\otimes B_j\). Its factors act on different bundle
factors and commute, so the product of these two skew-adjoint factors is
self-adjoint. Its smooth coefficients are bounded on \(M\), and its
Clifford factor makes it odd. Lemma 5.1a proves that their domains
remain the same \(H^1\) and that the compact resolvent is retained.
Lemma 5.1 makes their transforms compactly different, preserving the
compressed index.

Finally \(R_+=(1+D_E^2|_{H_+})^{-1/2}\) is an isomorphism from
\(L^2(S^+\otimes E)\) onto the domain of \(D_E^+\), with its graph norm.
The spectral-domain formula proves this exactly: the graph norm squared
is the weighted eigenvector sum of \(1+\lambda_j^2\) in the preceding
lesson, Lemma 4.0, and multiplication by \((1+\lambda^2)^{-1/2}\) gives
an isometry onto that domain, with inverse the reciprocal multiplier.
The estimate in the preceding lesson, Lemma 4.1b, identifies its graph
norm with \(H^1\). The equality \(F_{D_E}^+=D_E^+R_+\) therefore
identifies kernels and ranges with those of the Sobolev realization of
\(D_E^+\). Their indices agree. This proves (5.1). \(\square\)

The same argument shows independence from the connection. The
closed-manifold topological expression, with its orientation hypotheses
and symbol class, is proved in \href{KT-KK-15.html}{Dirac classes and
the cotangent Dolbeault element}, Sections 3--10. It is not a premise of
the analytic comparison proved here.

\subsection{6. Circle and noncommutative-torus
calculations}\label{circle-and-noncommutative-torus-calculations}

For the Hardy module \(h\) on the circle, multiplication by \(z^k\)
compresses to \(S^k\) for \(k\geq0\), or \((S^*)^{-k}\) for \(k<0\). The
former has zero kernel and cokernel dimension \(k\); the latter has
kernel dimension \(-k\) and zero cokernel. Thus \[
\langle[z^k],h\rangle=-k.
\tag{6.1}
\] For the matrix unitary \(\operatorname{diag}(z^2,z^{-5},1)\),
direct-sum additivity gives \[
-2+5+0=3.
\tag{6.2}
\] More generally the published prerequisite \emph{Toeplitz operators
and the index theorem on the circle}, Theorem 3.1, computes the matrix
Toeplitz index as minus the winding of its determinant. Formula (6.1)
already fixes the sign without using that more general calculation.

There is also an explicit Dirac-type cycle for a noncommutative torus.
Let \(A_\theta\) be the universal C*-algebra generated by unitaries
\(U,V\) with \(UV=e^{2\pi i\theta}VU\). On \(\ell^2(\mathbb Z^2)\), with
basis \(e_{m,n}\), represent them by \[
Ue_{m,n}=e_{m+1,n},\qquad
Ve_{m,n}=e^{-2\pi i\theta m}e_{m,n+1}.
\tag{6.3}
\] These are unitaries and satisfy the relation: on a basis vector the
two orders have phases differing by \(e^{2\pi i\theta}\). The universal
property therefore gives a representation. Its faithfulness is not
needed for the cycle.

On \(H=\ell^2(\mathbb Z^2)\otimes\mathbb C^2\), graded by \(\sigma_3\),
let \[
D(e_{m,n}\otimes v)
=e_{m,n}\otimes(m\sigma_1+n\sigma_2)v,
\tag{6.4}
\] with the square-summable graph domain from the preceding torus
example. That example's Fourier argument gives self-adjointness and
compact resolvent. Direct calculation gives \[
[D,U]=\sigma_1U,\qquad [D,V]=\sigma_2V.
\tag{6.5}
\] The shift bound \(1+(m\pm1)^2+n^2\leq3(1+m^2+n^2)\), and its
counterpart for shifting \(n\), prove domain preservation for
\(U,U^*,V,V^*\); the phase in \(V\) has modulus one and changes no graph
norm. The adjoint commutators are \([D,U^*]=-\sigma_1U^*\) and
\([D,V^*]=-\sigma_2V^*\). The product rule now proves domain
preservation and bounded commutators for every finite Laurent
polynomial. It is norm dense in \(A_\theta\) by its defining generation.
The bounded transform is therefore an even Fredholm module by the
preceding lesson, Theorem 4.1, for every real \(\theta\).

Testing the identity projection gives zero, since the two constant-mode
kernels each have dimension one, exactly as on the ordinary flat torus.
A smooth nontrivial projection \(p\) is tested by the typed compression
(1.1); this construction makes its index well defined without
presupposing a trace or a cyclic formula for that integer.

\subsection{7. Comparison with the Fredholm
character}\label{comparison-with-the-fredholm-character}

The compression indices admit trace formulas when their defects satisfy
a summability condition. We prove that comparison directly, using the
same kernel-minus-cokernel convention as in Sections 1--4.

Write \(\mathcal S_q(H_0,H_1)\), \(1\leq q<\infty\), for the compact
operators whose singular values have finite \(q\)-sum. The earlier
lesson
\href{../NCG-CYCLIC/public/reader/n-traces.html\#7-geometric-and-operator-examples}{\emph{Cyclic
forms that survive norm completion}, Lemma 7.3a} proves completeness,
finite-rank density, the bounded-multiplier estimate, inclusions
\(\mathcal S_q\subset\mathcal S_r\) for \(r\geq q\), and the full
Schatten Hölder inequality. In particular, \[
 \|X_1\cdots X_k\|_1\leq\prod_{j=1}^k\|X_j\|_k
 \tag{7.1}
\] for composable \(\mathcal S_k\) operators, also with bounded factors
inserted between them. Rectangular operators are included by putting
them in the appropriate off-diagonal corner of
\(\mathcal B(H_0\oplus H_1)\); their nonzero singular values and norms
are unchanged.

The ordinary trace is the trace of
\href{../foundations-of-von-neumann-algebras/compact-and-trace-class-operators-the-predual-of-b-h-and-the-operator-topologies.html\#OA-FND-LT-05}{\emph{Compact
and trace-class operators, the predual of \(B(H)\), and the operator
topologies}, Theorem 4.6}. That theorem proves basis independence,
continuity in trace norm, and cyclicity with a bounded multiplier.
Schatten products have cyclically invariant traces as well: approximate
each Schatten factor by finite-rank operators in its own norm, use (7.1)
and telescoping to pass every product to the trace-norm limit, and use
finite-dimensional cyclicity before taking the limit. For a trace-class
\(C:H_1\to H_0\) and bounded \(R:H_0\to H_1\), the same corner argument
gives \[
 \operatorname{Tr}_{H_0}(CR)=\operatorname{Tr}_{H_1}(RC).
 \tag{7.2}
\] These assertions never assign a trace to an individual factor which
is not trace class.

\subsubsection{7.1. The defect-power index
formula}\label{the-defect-power-index-formula}

\textbf{Lemma 7.1.} Let \(T:H_0\to H_1\) be bounded and suppose \[
 A_T=1_{H_0}-T^*T\in\mathcal S_q(H_0),\qquad
 B_T=1_{H_1}-TT^*\in\mathcal S_q(H_1).
 \tag{7.3}
\] Then \(T\) is Fredholm. For every integer \(m\geq q\), \[
 \operatorname{index}T
   =\operatorname{Tr}_{H_0}(A_T^m)
      -\operatorname{Tr}_{H_1}(B_T^m).
 \tag{7.4}
\]

\textbf{Proof.} First we obtain the closed range and the bounded inverse
needed in the trace calculation. On \(\ker T\), the compact operator
\(A_T\) acts as the identity. An infinite orthonormal family there would
have no convergent image subsequence, so \(\ker T\) is finite
dimensional. The same argument with \(B_T\) proves that \(\ker T^*\) is
finite dimensional.

There is a constant \(c>0\) with \(\|T\xi\|\geq c\|\xi\|\) on
\((\ker T)^\perp\). Otherwise choose unit vectors \(\xi_j\) in that
subspace with \(T\xi_j\to0\). Since \(T^*T\xi_j\to0\), compactness of
\(A_T\) supplies a subsequence on which \(\xi_j=A_T\xi_j+T^*T\xi_j\)
converges. Its limit is a unit vector in \((\ker T)^\perp\) annihilated
by \(T\), a contradiction. The lower bound proves that
\(\operatorname{ran}T\) is closed: a convergent sequence \(T\xi_j\) may
first have its kernel components removed, after which \(\xi_j\) is
Cauchy. The orthogonal complement of the range is \(\ker T^*\). Thus
\(T\) is Fredholm with the index of (0.1).

Let \(K\) and \(C\) be the orthogonal projections onto
\(\ker T\subset H_0\) and \(\ker T^*\subset H_1\). Define \[
 S_0:H_1\longrightarrow H_0
\] to be the inverse of \(T:(\ker T)^\perp\to\operatorname{ran}T\) on
that range and zero on its orthogonal complement. Its norm is at most
\(c^{-1}\). All operators here are bounded on their whole Hilbert
spaces; there is no unbounded-domain restriction. We have \[
 S_0T=1-K,\qquad TS_0=1-C,\qquad T^*C=0.
 \tag{7.5}
\] The typed difference \(R=T^*-S_0:H_1\to H_0\) satisfies \[
 R=-A_TS_0\in\mathcal S_q(H_1,H_0),
 \tag{7.6}
\] because \(A_TS_0=S_0-T^*TS_0=S_0-T^*(1-C)=S_0-T^*\).

Interpolate the inverse, keeping \(T\) fixed: \[
\begin{split}
 S_t&=S_0+tR,\\
 A_t&=1-S_tT=K-tRT,\\
 B_t&=1-TS_t=C-tTR,\qquad 0\leq t\leq1 .
\end{split}
 \tag{7.7}
\] The defects belong to \(\mathcal S_q\), and hence to
\(\mathcal S_m\), continuously and differentiably in those norms. They
need not be positive. They are self-adjoint in this construction, since
\(A_t=(1-t)K+tA_T\) and \(B_t=(1-t)C+tB_T\). Every \(m\)-fold product
used below is trace class by (7.1). The multiplication map from \(m\)
copies of \(\mathcal S_m\) to \(\mathcal S_1\) is continuous
multilinear. Expanding \(A_{t+h}^m\) gives the usual noncommutative
derivative, with trace-norm remainder bounded by a constant times
\(h^2\); each remainder term still has \(m\) Schatten factors. This also
covers \(m=1\), when the remainder is zero. Consequently, \[
\begin{split}
 \frac d{dt}\operatorname{Tr}(A_t^m)
   &=-m\operatorname{Tr}_{H_0}(RT A_t^{m-1}),\\
 \frac d{dt}\operatorname{Tr}(B_t^m)
   &=-m\operatorname{Tr}_{H_1}(TR B_t^{m-1}).
\end{split}
 \tag{7.8}
\] Cyclicity combines the \(m\) derivative summands; it is applied to
their complete trace-class products.

The exact identity \(TA_t=B_tT\) implies \(TA_t^{m-1}=B_t^{m-1}T\).
Moreover \(RB_t^{m-1}:H_1\to H_0\) is trace class: it has \(m\) factors
in \(\mathcal S_m\), including \(R\). Equations (7.2) and (7.8)
therefore give \[
 \operatorname{Tr}_{H_0}(RT A_t^{m-1})
   =\operatorname{Tr}_{H_0}(RB_t^{m-1}T)
   =\operatorname{Tr}_{H_1}(TR B_t^{m-1}).
\] The trace difference is constant. At \(t=0\), its two powers are the
finite-rank projections \(K,C\), so its value is
\(\dim\ker T-\dim\ker T^*\). At \(t=1\), \(S_1=T^*\) and the defects are
exactly (7.3). This proves (7.4). \(\square\)

\subsubsection{7.2. The odd compression and its
sign}\label{the-odd-compression-and-its-sign}

Use an exact self-adjoint involution \(F\), put \(P=(1+F)/2\), and let
\(U\) be the actual represented unitary of Section 2 or Section 3.
Amplifications are understood. Suppose \([F,U]\in\mathcal S_q\), and
choose an integer \(m\geq1\) with \(2m\geq q\). For \[
 T=PUP:PH\longrightarrow PH
\] its two defects, with identities on this closed range, are \[
\begin{split}
 A&=P-PU^*PUP=PU^*(1-P)UP,\\
 B&=P-PUPU^*P=PU(1-P)U^*P .
\end{split}
 \tag{7.9}
\] The operators \(V=(1-P)UP:PH\to(1-P)H\) and
\(W=PU(1-P):(1-P)H\to PH\) are corners of \([P,U]\), up to a sign. They
belong to \(\mathcal S_{2m}\). Therefore \(A=V^*V\) and \(B=WW^*\)
belong to \(\mathcal S_m\), and Lemma 7.1 gives \[
 \langle[u],x\rangle=\operatorname{Tr}_{PH}(A^m-B^m).
 \tag{7.10}
\]

There is an exact character version of this formula: \[
 \boxed{\quad
 \langle[u],x\rangle
  =\frac{(-1)^m}{2^{2m}}
     \operatorname{Tr}_H
       \left(F\bigl([F,U^*][F,U]\bigr)^m\right).
 \quad}
 \tag{7.11}
\] To prove it, write \(H=PH\oplus(1-P)H\). Since \([P,U^*]=-[P,U]^*\),
multiplication of these two off-diagonal blocks gives \[
 [F,U^*][F,U]=-4
       \begin{pmatrix}A&0\\0&W^*W\end{pmatrix}.
 \tag{7.12}
\] The traces of \((W^*W)^m\) and \((WW^*)^m=B^m\) are equal. One fully
typed verification approximates \(W:(1-P)H\to PH\) in
\(\mathcal S_{2m}\) by finite-rank operators. The \(2m\)-factor products
converge in trace norm by (7.1), and the two finite-rank traces are
equal before taking the limit. As \(F=1\oplus(-1)\), (7.12) and (7.10)
prove (7.11).

For the Hardy projection and the bilateral coordinate unitary \(U=z\),
\(T\) is the unilateral shift. Here \(A=0\) and \(B\) is the rank-one
projection onto the constant mode. Formula (7.10) gives \(-1\) in every
permitted degree. In degree one, (7.11) is specifically
\(-\frac14\operatorname{Tr}(F[F,U^*][F,U])\). This retains the positive
compression \(P=(1+F)/2\); replacing the compression or replacing \(U\)
by \(U^*\) changes the convention.

\subsubsection{7.3. The even compression}\label{the-even-compression-1}

Let \(\Gamma\) be the grading, let \(F\) be odd and satisfy \(F=F^*\),
\(F^2=1\), and write \(e=\rho^{(n)}(p)\) for a represented projection.
Put \(e_\pm=\rho_\pm^{(n)}(p)\) and retain the typed operator of (1.1),
\[
 T_p=e_-F_+e_+:e_+H_+^n\longrightarrow e_-H_-^n .
\] Assume \([F,e]\in\mathcal S_q\) and choose \(m\geq1\) with
\(2m\geq q\). Exactness makes \(F_-=F_+^*\), so the inverse errors are
\[
\begin{split}
 A_p&=e_+-T_p^*T_p
      =e_+F_-(1-e_-)F_+e_+,\\
 B_p&=e_--T_pT_p^*
      =e_-F_+(1-e_+)F_-e_- .
\end{split}
 \tag{7.13}
\] Each is a product of two \(\mathcal S_{2m}\) corners of \([F,e]\),
with bounded factors. Thus they lie in \(\mathcal S_m\) on their
respective, possibly different, Hilbert spaces. Lemma 7.1 proves \[
 \langle[p],x\rangle
    =\operatorname{Tr}_{e_+H_+^n}(A_p^m)
       -\operatorname{Tr}_{e_-H_-^n}(B_p^m).
 \tag{7.14}
\]

Writing \(D_e=[F,e]\), its two diagonal corners relative to \(e\) are
zero, so \(D_e^2\) commutes with \(e\). Direct multiplication gives \[
 eD_e^2e=-eF(1-e)Fe
         =-\begin{pmatrix}A_p&0\\0&B_p\end{pmatrix}
             \quad\hbox{on }eH^n.
\] Taking powers and the grading trace in (7.14) gives \[
 \langle[p],x\rangle
       =(-1)^m\operatorname{Tr}_{H^n}
                         \bigl(\Gamma e[F,e]^{2m}\bigr).
 \tag{7.15}
\] Its full bounded-character form is \[
 \boxed{\quad
 \langle[p],x\rangle
   =\frac{(-1)^m}{2}\operatorname{Tr}_{H^n}
                         \bigl(\Gamma F[F,e]^{2m+1}\bigr).
 \quad}
 \tag{7.16}
\] Here every displayed trace is defined, since \(2m\) commutators
already have a trace-class product. For the identity between the two
traces, set \(\omega=e[F,e]^{2m}\). On operator forms use the derivation
\[
 d\omega=F\omega-(-1)^{\deg\omega}\omega F .
 \tag{7.17}
\] It has \(d^2=0\), since \(F^2=1\), and \(d\omega=[F,e]^{2m+1}\).
Cyclicity and \(F\Gamma F=-\Gamma\) give \[
 \tfrac12\operatorname{Tr}(\Gamma Fd\omega)
   =\tfrac12\bigl(\operatorname{Tr}(\Gamma\omega)
                  -\operatorname{Tr}(\Gamma F\omega F)\bigr)
   =\operatorname{Tr}(\Gamma\omega).
\] This proves (7.16) with the positive-to-negative block, including its
sign.

\subsubsection{7.4. The character cocycles and their
normalization}\label{the-character-cocycles-and-their-normalization}

Suppose that an exact representative has a norm-dense \(*\)-subalgebra
\(\mathcal D\subset A\) with \([F,\rho(a)]\in\mathcal S_q\) for
\(a\in\mathcal D\). Coefficients act through the unital extension
\(\rho^+\) on the forced unitization; this is also the unitary
completion of Section 2 in the unital case. Suppress \(\rho^+\) in the
products. The scalar matrix trace is included after amplification.

For a positive even degree \(n=2m\geq q\), and an odd degree
\(n=2m-1\geq q\), respectively, define the raw bounded Fredholm
characters \[
\begin{split}
 \kappa_{2m}(a_0,\ldots,a_{2m})
  &=\tfrac12\operatorname{Tr}
             \bigl(\Gamma F[F,a_0]\cdots[F,a_{2m}]\bigr),\\
 \kappa_{2m-1}(a_0,\ldots,a_{2m-1})
  &=\tfrac12\operatorname{Tr}
             \bigl(F[F,a_0]\cdots[F,a_{2m-1}]\bigr).
\end{split}
 \tag{7.18}
\] The condition \(n\geq q\) is sufficient for the cocycle and
coefficient estimates below; the individual compression identities
(7.11) and (7.16) only need the weaker indicated \(2m\geq q\).

Here is a proof of their cyclic-cocycle assertions. Represent a
degree-\(k\) operator form by \(a_0da_1\cdots da_k\), retaining the
degree as part of the form. The derivation (7.17) agrees with
\(da=[F,a]\), obeys the graded Leibniz rule, and has square zero. Forms
of degree \(n\) are trace class by (7.1). In even degree use
\(\int\omega=\operatorname{Tr}(\Gamma\omega)\), and in odd degree use
\(\int\omega=\operatorname{Tr}(\omega)\).

These are graded traces in degree \(n\). In the even case, moving a
degree-\(r\) factor past \(\Gamma\) gives \((-1)^r=(-1)^{r(n-r)}\). In
the odd case \(r(n-r)\) is even, so ordinary cyclicity has precisely the
required graded sign. All these rotations are justified for their
complete Schatten products. In both cases, \[
 \int\omega=\tfrac12\operatorname{Tr}(J Fd\omega),
 \qquad
 J=\Gamma\ \hbox{in even degree},\quad J=1\ \hbox{in odd degree}.
 \tag{7.19}
\] The even identity was proved above. For odd \(\omega\), it follows
from \(d\omega=F\omega+\omega F\) and
\(\operatorname{Tr}(F\omega F)=\operatorname{Tr}(\omega)\). Thus
\(\int d\eta=0\), by \(d^2=0\), without assigning a trace to a
degree-\((n-1)\) form \(\eta\). Equation (7.19) also proves
\(\kappa_n(a_0,\ldots,a_n)=\int a_0da_1\cdots da_n\).

Expand the Hochschild differential using \(d(ab)=(da)b+a\,db\).
Consecutive terms cancel, leaving \[
 (-1)^n\int a_0da_1\cdots da_n\,a_{n+1}
   +(-1)^{n+1}\int a_{n+1}a_0da_1\cdots da_n=0 .
\] Hence \(b\kappa_n=0\). Closedness applied to
\(a_na_0da_1\cdots da_{n-1}\), followed by moving \(da_n\) around the
graded trace, gives \[
 \kappa_n(a_n,a_0,\ldots,a_{n-1})
         =(-1)^n\kappa_n(a_0,\ldots,a_n).
\] This proves cyclicity. It also gives the coefficient estimate \[
 \left|\int x_1da_1\cdots x_nda_n\right|
     \leq\prod_{j=1}^n\|\rho^+(x_j)\|
                     \prod_{j=1}^n\|[F,\rho^+(a_j)]\|_n ,
 \tag{7.20}
\] again by (7.1). Thus the characters have the \(n\)-trace estimate
with the actual unitized coefficient action.

Fix the cyclic evaluation convention \[
\begin{aligned}
 J_\tau[p]&=\tau^{(N)}(p,\ldots,p)
                    &&\text{in even degree},\\
 J_\tau[u]&=\tau^{(N)}(u^*,u,\ldots,u^*,u)
                    &&\text{in odd degree}.
\end{aligned}
 \tag{7.21}
\] This is the convention of
\href{../NCG-CYCLIC/public/reader/n-traces.html\#4-the-form-reaches-k-theory}{\emph{Cyclic
forms that survive norm completion}, Theorem 4.3}; it puts no additional
factorial in these evaluations. With that explicit convention the
index-normalized characters are \[
 \operatorname{Ch}_{2m}=(-1)^m\kappa_{2m},\qquad
 \operatorname{Ch}_{2m-1}
       =\frac{(-1)^m}{2^{\,2m-1}}\kappa_{2m-1}.
 \tag{7.22}
\] Equations (7.11) and (7.16) now state exactly \[
 J_{\operatorname{Ch}_{2m}}[p]=\langle[p],x\rangle,\qquad
 J_{\operatorname{Ch}_{2m-1}}[u]=\langle[u],x\rangle .
 \tag{7.23}
\] All entries and traces use the amplified representation, including
for a rectangular even compression.

\subsubsection{7.5. Relative classes and all K-theory
representatives}\label{relative-classes-and-all-k-theory-representatives}

For a nonunital algebra use the exact representative and
forced-unitization action of Theorem 3.1. We require the stated Schatten
condition on that representative; exact normalization of a general
compact-defect cycle alone is not a Schatten estimate. If \([p]-[r]\) is
a relative class with equal scalar ranks, (7.23) gives \[
 \langle[p]-[r],x\rangle
    =J_{\operatorname{Ch}_{2m}}[p]
       -J_{\operatorname{Ch}_{2m}}[r].
 \tag{7.24}
\] A scalar projection commutes exactly with \(F\), so its
positive-degree character is zero; its compression is a unitary between
the two scalar-selected graded ranges. A scalar unitary also commutes
with \(P\), has unitary compression, and has zero character. The old
identity of a unital \(A\), viewed in its forced unitization, is not the
newly adjoined scalar identity. Only the latter is subject to this
scalar-zero assertion. This preserves the ordinary unit pairing of
Section 1.

There is no additional assumption that every ambient K-class already has
a representative in an arbitrarily chosen smaller smooth algebra. For
the Schatten domain itself put \[
 \mathcal D_F=\{a\in A:[F,\rho(a)]\in\mathcal S_q\},\qquad
 \|a\|_{\mathcal D_F}=\|a\|+\|[F,\rho(a)]\|_q .
\] It is a \(*\)-algebra: the commutator product rule and the ideal
estimate give closure under products, and
\([F,\rho(a^*)]=-[F,\rho(a)]^*\) gives closure under adjoints.
Contractivity of \(\rho\) also gives
\(\|ab\|_{\mathcal D_F}\leq\|a\|_{\mathcal D_F}\|b\|_{\mathcal D_F}\).
Its derivation is closed: convergence in \(\mathcal S_q\) implies
operator-norm convergence, so if \(a_j\to a\) in \(A\) and
\([F,\rho(a_j)]\to R\) in \(\mathcal S_q\), then \(R=[F,\rho(a)]\).
Completeness of \(A\) and \(\mathcal S_q\) proves completeness in the
displayed graph norm. The same statements hold in every finite matrix
size. We use the continuous functional calculus and its norm continuity
proved in
\href{../foundations-of-von-neumann-algebras/c-star-algebras-continuous-functional-calculus-automatic-continuity-positive-cones.html\#OA-FND-CF-07}{\emph{C*-algebras:
continuous functional calculus, automatic continuity, positive cones,
approximate identities and quotients}, Theorems 5.1 and 6.1}.

For an invertible matrix \(a\) over \(A^+\) lying in this domain, \[
 [F,\rho^+(a^{-1})]
     =-\rho^+(a^{-1})[F,\rho^+(a)]\rho^+(a^{-1}).
 \tag{7.25}
\] Consequently the domain is closed under resolvents. If a self-adjoint
matrix \(b\) in the domain is sufficiently close to a projection \(p\),
its spectrum lies in disjoint small neighborhoods of \(0,1\). The
projection \[
 e=\frac1{2\pi i}\int_\gamma(z-b)^{-1}\,dz
\] for a positively oriented circle around the neighborhood of \(1\)
lies in the domain. Indeed the resolvents are continuous in operator
norm and their commutators are continuous in \(\mathcal S_q\) by (7.25),
so both contour integrals exist in their respective Banach norms.
Continuous functional calculus for the self-adjoint \(b\) identifies
this integral with its spectral cut: the scalar circle integral is one
inside and zero outside, by the geometric series on a circle. The
Neumann resolvent identity also shows that \(e\to p\) as \(b\to p\).
Thus they are equivalent by the close-projection unitary of the
preceding lesson. A relative projection can be approximated while fixing
its scalar part, which this cut preserves.

Likewise approximate a unitary \(u\) by a matrix \(v\) in the domain,
sufficiently closely to make \(v\) invertible. Its polar unitary
\(w=v(v^*v)^{-1/2}\) belongs to the domain. To check the inverse square
root without a differentiability assumption on functional calculus,
choose \(c>0\) with \(\|1-cv^*v\|<1\). The binomial series \[
 (v^*v)^{-1/2}
    =\sqrt c\sum_{k\geq0}\frac{\binom{2k}{k}}{4^k}(1-cv^*v)^k
\] converges in the graph norm: its commutator is bounded term by term
by \(k\|1-cv^*v\|^{k-1}\|[F,1-c\rho^+(v^*v)]\|_q\), and the resulting
scalar series converges. If \(v\to u\), the binomial estimate or
continuous functional calculus gives \(w\to u\). Close unitaries are
homotopic, for example by polarizing their straight segment, which
remains invertible. Approximations fixing the scalar unit of a relative
unitary give \(w\) with that same scalar unit.

When \(\mathcal D_F\) is norm dense, these arguments supply domain
representatives for every K-class. Formulas (7.23)--(7.24) on those
representatives are the already-defined analytic pairings. Their
independence, bilinearity, homotopy invariance and relative scalar
convention therefore follow from Sections 1--3, with no additional
summability requirement on a path used to identify ambient K-classes.
This proves the bounded Fredholm-character comparison on the full
K-theory pairing.

The equality with Kasparov products is proved for even projection
classes by \href{KT-KK-09.html}{Connections and the existence of the
Kasparov product}, Solution to Exercise 10.2, and for odd classes by
\href{KT-KK-10.html}{Homotopy, associativity, the index pairing and
KK-equivalence}, Theorem 6.2, with the ordered Clifford and
positive-exponential conventions there. Under
\(K_i(A)=KK_i(\mathbb C,A)\) and \(K^i(A)=KK_i(A,\mathbb C)\), the
product has the compression index of Sections 1--4 in
\(KK_0(\mathbb C,\mathbb C)=\mathbb Z\), with those signs. The trace
proof here does not require that product construction. Residue formulas
for unbounded spectral triples require further analytic hypotheses and a
separate comparison theorem; none is a premise of (7.4)--(7.24).

\subsection{8. Exercises}\label{exercises}

\textbf{Exercise 8.1 (basic).} Compute \(\langle[z^k],h\rangle\) for
every integer \(k\), including zero. Compute the pairing of
\(\operatorname{diag}(z^2,z^{-5},1)\).

\textbf{Exercise 8.2 (intermediate).} Prove that a compact perturbation
of the odd cycle operator preserves the pairing. Explain why one cannot
simply call \((1+F_t)/2\) a projection along an arbitrary perturbation
segment.

\textbf{Exercise 8.3 (intermediate).} Prove
\(\langle[uv],x\rangle=\langle[u],x\rangle+\langle[v],x\rangle\) for
unitaries of a common matrix size. Include the stable-unitary group law.

\textbf{Exercise 8.4 (advanced).} Use an essentially unitary Fredholm
lift \(T\) to prove that the odd pairing equals the extension boundary.
Keep track of both defect projections and evaluate the result for the
unilateral shift.

\subsection{9. Solutions}\label{solutions}

\textbf{Solution 8.1.} For \(k>0\), \(S^k\) has cokernel spanned by
\(e_0,\ldots,e_{k-1}\) and zero kernel, giving \(-k\). For
\(k=-\ell<0\), \((S^*)^\ell\) is surjective with kernel spanned by
\(e_0,\ldots,e_{\ell-1}\), giving \(\ell=-k\). For \(k=0\) the
compression is the identity, giving zero. The three diagonal entries
contribute \(-2,+5,0\), so the matrix pairing is \(3\).

\textbf{Solution 8.2.} Let \(F'=F+K\) with \(K\) compact and both
endpoints cycles. The segment \(F_t=F+tK\) is a cycle homotopy by the
explicit defect expansion in the preceding lesson, Solution 8.2. Its
exact normalizations \(N(F_t)\) form a norm-continuous path of
self-adjoint involutions. Their positive projections \(P_t\) therefore
form a norm-continuous projection path. The close-projection unitaries
in that lesson's Lemma 7.1 transport their ranges locally to a fixed
space, giving norm-continuous Fredholm compressions, with constant
index. This proves the claim. Without normalization the defect \[
\left(\frac{1+F_t}{2}\right)^2-\frac{1+F_t}{2}
=\frac{F_t^2-1}{4}
\] need not be zero. Compactness, or local compactness, of this defect
does not make the operator an actual projection.

\textbf{Solution 8.3.} For exact normalization let \(U,V\) be the
represented unitary completions. Compactness of \([P,V]\) gives \[
PUVP-(PUP)(PVP)=PU(1-P)VP\in\mathcal K(PH).
\] The compressions \(PUP\), \(PVP\) are Fredholm. Compact invariance
and composition additivity make the index of the left product equal the
sum of their indices. The rotation path in Theorem 2.2 identifies the
stable block sum of \(u,v\) with \(uv\) plus an identity block, so this
index sum is the addition in \(K_1(A)\), not a different operation on
representatives. The unitization formulas of Theorem 3.1 give the same
proof for a nonunital algebra.

\textbf{Solution 8.4.} Fredholm closed range gives a spectral gap for
\(|T|\) off its finite-dimensional kernel. Define \(h(0)=0\),
\(h(t)=1/t\) on the nonzero spectrum, and \(V=Th(|T|)\). This is the
polar partial isometry, with \(V^*V=1-p_k\), \(VV^*=1-p_c\). Essential
unitarity gives \(|T|-1\) compact and \(T-V=V(|T|-1)\) compact. Thus
\(V\) lifts the same quotient unitary. The matrix \(W\) in (4.2) is
unitary: its diagonal products are \(VV^*+p_c=1\), \(p_k+V^*V=1\), and
its off-diagonal products vanish because \(Vp_k=p_cV=0\). It is a
doubled lift, and (4.3) gives boundary \([p_k]-[p_c]\). The rank
identification gives the Fredholm index. For \(T=S\), the kernel
projection is zero and the cokernel projection is \(p_0\), so the
boundary is \(-[p_0]\) and the integer pairing is \(-1\).

\subsection{Internal proof dependencies and source
credit}\label{internal-proof-dependencies-and-source-credit}

The preceding lesson, Lemma 2.1 and Propositions 2.2--2.3, proves the
locally compact ideal and continuous normalization; its Lemma 7.1 proves
range transport by explicit close-projection unitaries. Lemma 1.1 here
gives all typed even inverse and adjoint errors, Lemma 2.1 gives the odd
errors and product rule, Lemma 4.0 proves the full relative compact-rank
identification, and Lemmas 5.1--5.1a give the compact transform
comparison and domain preservation used for twisting. The operator
K-theory lessons \emph{Nonunital algebras: unitization, relative classes
and half-exactness} and \emph{Invertibles, unitaries and \(K_1\)} supply
the exact relative and stable-unitary definitions; \emph{The index map
and the exact sequence at \(K_0\)}, Definition 2.1 and Theorem 2.2,
supplies the doubled-lift boundary and its naturality. \emph{Toeplitz
operators and the index theorem on the circle}, Theorem 3.1, supplies
the general matrix-symbol calculation. The Hilbert-space Fredholm
prerequisites are those identified in the preceding lesson. Its local
Lemmas 4.1a--4.1c and Corollary 4.2 prove the entire closed-manifold
Sobolev, elliptic-estimate and adjoint-domain foundation used by Theorem
5.2 here.

The bounded Fredholm-character comparison is proved in Section 7 using
the exact Schatten and ordinary-trace foundations named there. The
Kasparov-product interpretation is proved in the earlier product lessons
identified in Section 7. The closed-manifold topological Dirac index
comparison is proved in \href{KT-KK-15.html}{Dirac classes and the
cotangent Dolbeault element}, Sections 3--10. The pairings, their
invariance, the relative scalar check, the extension boundary sign, the
analytic twisting bridge, and all examples and exercises above are
proved here.

\subsection{References}\label{references}

\begin{itemize}
\item
  B. Blackadar, \emph{K-Theory for Operator Algebras}, freely accessible
  author-corrected second-edition PDF. Sections 16.3.2, 17.5 and 18.10
  develop the Fredholm, boundary and index-pairing pictures.
  \href{https://www.bruceblackadar.com/Mathematics/book6.pdf}{Free
  author version}. The typed compression, relative scalar check, doubled
  boundary lift, compact-rank identification and analytic twisting
  argument required here are proved in this lesson.
\item
  J. Rosenberg, \emph{Examples and Applications of Noncommutative
  Geometry and K-Theory}, author notes dated 2010, Sections 1.2 and 4.2.
  \href{https://math.umd.edu/~jmr/BuenosAires/NCGexamples.pdf}{Free
  author notes}. These sections supply the source comparison for
  Fredholm cycles and smooth noncommutative tori; the exact
  representation, graph domain and commutators of Section 6 are verified
  here.
\end{itemize}

These free source PDFs retain their authors' copyrights. The CC0 notice
applies to this lesson's original text.

\end{document}
